\chapter{Agent Architectures and Topologies}\label{ch:architectures}

Agent orchestration topologies represent the fundamental structural patterns through which multiple agents coordinate to solve complex problems. Each topology embodies distinct trade-offs between latency, parallelism, complexity, and observability. Understanding these patterns is essential for selecting or designing an orchestration architecture that matches problem domain constraints, operational budgets, and quality requirements~\cite{LangChain2024a}.

\section{Linear Sequential Orchestration}\label{sec:linear_seq}

Linear sequential orchestration chains agents in strict order: Agent A completes its task, produces structured output, passes it to Agent B, which processes and passes results to Agent C, and so forth. This pattern mirrors traditional data pipelines and is ideal for problems with a natural progression of steps. A research pipeline exemplifies this topology: a Researcher agent gathers information and synthesizes sources, a Writer agent drafts narrative content from the synthesis, an Editor agent refines prose and enforces style guidelines, and a \LaTeX{} Producer agent formats the final document for publication~\cite{CrewAI2024}.

\begin{figure}[H]\centering\resizebox{\textwidth}{!}{% TikZ block diagram of the CrewAI sequential pipeline
% Included via % TikZ block diagram of the CrewAI sequential pipeline
% Included via % TikZ block diagram of the CrewAI sequential pipeline
% Included via \input{figures/crew_architecture.tikz}
\begin{tikzpicture}[
  node distance=1.6cm and 2.4cm,
  agent/.style={rectangle, rounded corners=4pt, draw=uohblue, fill=uohblue!10,
                minimum width=2.8cm, minimum height=1.0cm, align=center,
                font=\small\bfseries},
  tool/.style={rectangle, draw=gray!70, fill=gray!10,
               minimum width=2.2cm, minimum height=0.7cm, align=center,
               font=\scriptsize},
  arrow/.style={-Latex, thick, color=uohblue},
  darrow/.style={-Latex, dashed, color=gray},
]

% Agents
\node[agent] (researcher) {Researcher\\Agent};
\node[agent, right=of researcher] (writer) {Writer\\Agent};
\node[agent, right=of writer] (editor) {Editor\\Agent};
\node[agent, right=of editor] (latex) {LaTeX\\Agent};

% Pipeline arrows
\draw[arrow] (researcher) -- (writer);
\draw[arrow] (writer) -- (editor);
\draw[arrow] (editor) -- node[above, font=\tiny, color=gray]{approve?} (latex);

% Tools below each agent
\node[tool, below=0.8cm of researcher] (t1) {WebSearchTool\\FileWriteTool};
\node[tool, below=0.8cm of writer]    (t2) {FileReadTool\\FileWriteTool};
\node[tool, below=0.8cm of editor]    (t3) {FileReadTool\\FileWriteTool};
\node[tool, below=0.8cm of latex]     (t4) {ChartGen\\LaTeXCompile};

\draw[darrow] (researcher) -- (t1);
\draw[darrow] (writer) -- (t2);
\draw[darrow] (editor) -- (t3);
\draw[darrow] (latex) -- (t4);

% SDK + Gatekeeper box
\begin{scope}[on background layer]
  \node[draw=orange!60, fill=orange!5, rounded corners, inner sep=8pt,
        fit=(researcher)(writer)(editor)(latex)(t1)(t4),
        label={[font=\scriptsize, color=orange!80]above:ArticleSDK via ApiGatekeeper}] {};
\end{scope}

\end{tikzpicture}

\begin{tikzpicture}[
  node distance=1.6cm and 2.4cm,
  agent/.style={rectangle, rounded corners=4pt, draw=uohblue, fill=uohblue!10,
                minimum width=2.8cm, minimum height=1.0cm, align=center,
                font=\small\bfseries},
  tool/.style={rectangle, draw=gray!70, fill=gray!10,
               minimum width=2.2cm, minimum height=0.7cm, align=center,
               font=\scriptsize},
  arrow/.style={-Latex, thick, color=uohblue},
  darrow/.style={-Latex, dashed, color=gray},
]

% Agents
\node[agent] (researcher) {Researcher\\Agent};
\node[agent, right=of researcher] (writer) {Writer\\Agent};
\node[agent, right=of writer] (editor) {Editor\\Agent};
\node[agent, right=of editor] (latex) {LaTeX\\Agent};

% Pipeline arrows
\draw[arrow] (researcher) -- (writer);
\draw[arrow] (writer) -- (editor);
\draw[arrow] (editor) -- node[above, font=\tiny, color=gray]{approve?} (latex);

% Tools below each agent
\node[tool, below=0.8cm of researcher] (t1) {WebSearchTool\\FileWriteTool};
\node[tool, below=0.8cm of writer]    (t2) {FileReadTool\\FileWriteTool};
\node[tool, below=0.8cm of editor]    (t3) {FileReadTool\\FileWriteTool};
\node[tool, below=0.8cm of latex]     (t4) {ChartGen\\LaTeXCompile};

\draw[darrow] (researcher) -- (t1);
\draw[darrow] (writer) -- (t2);
\draw[darrow] (editor) -- (t3);
\draw[darrow] (latex) -- (t4);

% SDK + Gatekeeper box
\begin{scope}[on background layer]
  \node[draw=orange!60, fill=orange!5, rounded corners, inner sep=8pt,
        fit=(researcher)(writer)(editor)(latex)(t1)(t4),
        label={[font=\scriptsize, color=orange!80]above:ArticleSDK via ApiGatekeeper}] {};
\end{scope}

\end{tikzpicture}

\begin{tikzpicture}[
  node distance=1.6cm and 2.4cm,
  agent/.style={rectangle, rounded corners=4pt, draw=uohblue, fill=uohblue!10,
                minimum width=2.8cm, minimum height=1.0cm, align=center,
                font=\small\bfseries},
  tool/.style={rectangle, draw=gray!70, fill=gray!10,
               minimum width=2.2cm, minimum height=0.7cm, align=center,
               font=\scriptsize},
  arrow/.style={-Latex, thick, color=uohblue},
  darrow/.style={-Latex, dashed, color=gray},
]

% Agents
\node[agent] (researcher) {Researcher\\Agent};
\node[agent, right=of researcher] (writer) {Writer\\Agent};
\node[agent, right=of writer] (editor) {Editor\\Agent};
\node[agent, right=of editor] (latex) {LaTeX\\Agent};

% Pipeline arrows
\draw[arrow] (researcher) -- (writer);
\draw[arrow] (writer) -- (editor);
\draw[arrow] (editor) -- node[above, font=\tiny, color=gray]{approve?} (latex);

% Tools below each agent
\node[tool, below=0.8cm of researcher] (t1) {WebSearchTool\\FileWriteTool};
\node[tool, below=0.8cm of writer]    (t2) {FileReadTool\\FileWriteTool};
\node[tool, below=0.8cm of editor]    (t3) {FileReadTool\\FileWriteTool};
\node[tool, below=0.8cm of latex]     (t4) {ChartGen\\LaTeXCompile};

\draw[darrow] (researcher) -- (t1);
\draw[darrow] (writer) -- (t2);
\draw[darrow] (editor) -- (t3);
\draw[darrow] (latex) -- (t4);

% SDK + Gatekeeper box
\begin{scope}[on background layer]
  \node[draw=orange!60, fill=orange!5, rounded corners, inner sep=8pt,
        fit=(researcher)(writer)(editor)(latex)(t1)(t4),
        label={[font=\scriptsize, color=orange!80]above:ArticleSDK via ApiGatekeeper}] {};
\end{scope}

\end{tikzpicture}
}\caption{CrewAI multi-agent architecture.}\label{fig:crew_architecture}\end{figure}

Sequential architectures offer predictability and simplicity: you know exactly when each agent runs, what input it receives, and in what order results are computed. The trade-off is rigidity. If an early agent produces unusable output, the entire pipeline must restart or be manually corrected. Additionally, latency accumulates linearly with the number of agents; there are no parallelism opportunities to hide latency under concurrent execution~\cite{LangChain2024a}.

The latency of a sequential pipeline is given by:
\begin{equation}L_{\mathrm{seq}}=\sum_{i=1}^{N}\ell_i\label{eq:pipeline_latency}\end{equation}
where $\ell_i$ is the latency (wall-clock time) of agent $i$ and $N$ is the total number of agents. For a 4-agent research pipeline with per-agent latencies of 5\,s (Researcher), 10\,s (Writer), 8\,s (Editor), and 6\,s (\LaTeX{} Producer), the total pipeline latency is 29\,s, unavoidable by design.

\section{Hierarchical Manager-Worker Orchestration}\label{sec:hierarchical}

Hierarchical systems introduce a supervisor agent (manager) that observes problem state and delegates tasks to worker agents based on dynamic conditions. The supervisor reasons about the problem structure, selects which workers to invoke, synthesizes their outputs, and decides whether additional work is needed. This pattern enables adaptive execution: the supervisor can route complex queries to specialist agents and simple queries to faster, cheaper models, achieving latency and cost optimization not possible in rigid sequential pipelines~\cite{LangChain2024a}.

Hierarchical topologies scale well to many specialized agents (10+) because the supervisor provides a single decision point, preventing exponential branching of execution paths. However, the cost is complexity: the supervisor itself must be sophisticated enough to route reliably, and debugging routing errors requires introspection of the supervisor's reasoning process. Additionally, the supervisor becomes a bottleneck if it cannot parallelize its worker invocations.

\section{Peer-to-Peer Consensus Orchestration}\label{sec:peer_consensus}

Peer-to-peer topologies allow agents to work in parallel without a central coordinator, communicating directly with one another to reach consensus on shared decisions. Multiple researchers might simultaneously investigate different aspects of a topic, then hand off to a shared editor who synthesizes and cross-validates their findings. Parallelism reduces wall-clock latency and enables true concurrency; a 4-agent peer team can complete work in approximately $\max_i(\ell_i)$ time rather than $\sum_i(\ell_i)$~\cite{AutoGen2023}.

However, peer-to-peer systems require careful state management: shared resources must be protected from concurrent modification, and convergence points must handle variable-speed agents gracefully. Debugging is harder because there is no central audit trail; execution traces must be reconstructed from each agent's logs. Additionally, consensus mechanisms (voting, averaging, arbitration) add latency and complexity.

\section{Tool-Augmented Specialist Orchestration}\label{sec:tool_augmented}

Tool-augmented orchestration provides agents with rich toolkits (web search, code execution, file I\texttt{\&}O, external API access) rather than expecting them to reason from static knowledge. A Researcher agent equipped with web search tools can fetch the latest information; a Writer agent equipped with a chart-generation tool can create visualizations programmatically; a \LaTeX{} Producer agent equipped with a compilation tool can validate syntax and produce PDFs. This pattern significantly improves output quality and grounding in observed reality rather than hallucination~\cite{LangChain2024a}.

Tool access introduces new operational concerns: tools are rate-limited (web search APIs allow $N$ calls per minute), carry cost (each tool invocation consumes tokens and API credits), and may fail (network errors, malformed outputs). Production systems implement a Gatekeeper layer that intercepts all tool calls, enforces rate limits, tracks token consumption, and handles errors gracefully~\cite{CrewAI2024}.

\section{Topology Comparison and Selection}\label{sec:topology_comparison}

Each topology presents distinct trade-offs. The following table synthesizes the key dimensions:

\begin{longtable}{@{}p{2.8cm}p{2.8cm}p{2.6cm}p{2.0cm}p{3.9cm}@{}}
\toprule
\textbf{Topology} & \textbf{Latency Model} & \textbf{Parallelism} & \textbf{Complexity} & \textbf{Ideal For} \\[5pt]
\endfirsthead
\toprule
\textbf{Topology} & \textbf{Latency Model} & \textbf{Parallelism} & \textbf{Complexity} & \textbf{Ideal For} \\[5pt]
\endhead
\midrule
Linear Sequential & $\sum_i \ell_i$ & None & Low & Well-defined workflows: research pipelines, content generation, refining stages \\[5pt]
Hierarchical & $\ell_{\mathrm{super}} + \max_i(\ell_i^{\mathrm{worker}})$ & Within supervisor's workers & Medium & Adaptive routing: complex tasks routed to specialists, simple to fast models \\[5pt]
Peer-to-Peer & $\max_i(\ell_i) + \ell_{\mathrm{sync}}$ & Full & High & Consensus problems: debate, fact-checking, multiple independent investigations \\[5pt]
Tool-Augmented & Depends on tool costs & If tools async & Medium & Grounding in reality: web search, code execution, external data integration \\[5pt]
Graph-Based (DAG) & Depends on structure & Full, conditional & High & General workflows: cycles, branches, complex dependencies \\[5pt]
\bottomrule
\end{longtable}

Selection guidance follows from problem characteristics:

\begin{itemize}
\item \textbf{Sequential}: Choose when the problem has a natural linear progression (research $\to$ draft $\to$ edit $\to$ format). Pros: simple to implement, easy to debug. Cons: no parallelism, bottleneck on slow stages.
\item \textbf{Hierarchical}: Choose when many specialized agents must be coordinated by a decision-maker. Pros: adaptive, scalable to many agents. Cons: supervisor complexity, potential bottleneck.
\item \textbf{Peer-to-Peer}: Choose when multiple agents must independently verify a claim or explore a problem from different angles. Pros: parallelism, robustness. Cons: consensus complexity, harder debugging.
\item \textbf{Tool-Augmented}: Choose when agents must ground outputs in real-world data (live search, APIs, computations). Pros: better quality, less hallucination. Cons: operational burden (rate limits, errors, cost).
\item \textbf{Graph-Based (DAG)}: Choose when the workflow has complex branching, loops, or conditional dependencies. Pros: maximally expressive. Cons: complex to design and debug.
\end{itemize}

Production deployments often combine topologies: use sequential stages within a hierarchical supervisor, or use peer consensus within a graph-based framework to handle conditional branches. The key insight is that topology selection shapes latency, cost, and failure modes; misaligned topology and problem structure lead to poor performance~\cite{Anthropic2024b}.